\chapter{Diagnosis of the In-House File Decision}
\label{app:diagnosis}
% 2026-11-26: 7장 VII-3~VII-5(원래 본문)를 옮겼다. 절 label 은 그대로라 본문의 참조가 이 부록을 가리킨다.

\section{A recorded file-processing decision}
\label{sec:recorded-case}
% 2026-11-26 심사(Sonnet P2, Opus P2-5): 4장 IV-4 를 옮겼다.

Two terms of the synthetic procedure are needed here. A \emph{payload
search} is the synthetic file search (DISCOVERY) that looks for a
target's payload file, and its exit status records whether it found
one. A \emph{marker search} is a later search on the same target that
marks it as the current one; its own exit status is not the payload
result. Execution \texttt{THME0064} comes from an earlier three-target
collection of this project (Appendix~\ref{app:secondary-results}), not
from the C-fold data. Table~\ref{tab:recorded-case} shows part of its
received prefix before its first file-processing decision
(recorded as \texttt{SEQ\_5}). The path
names distinguish three synthetic targets. The marker concerns
\texttt{alpha}, whose earlier payload search returned a nonzero exit
status, while the marker's own zero exit status belongs to a different
process. The most recent exit status received in the \texttt{alpha}
group is therefore the marker's zero, even though the payload search
on that target failed; a model that read only the latest status would
see a success that the search never had. Grouping the records by
target lets a model read the two processes together while keeping
their results apart. In this execution the action that followed was
again a synthetic file search. The example illustrates the kind of
target-result lookup examined in this appendix.

\begin{table}[htbp]
\centering
\caption[Selected records received before the decision in \texttt{THME0064}]{Selected records received before the decision in
\texttt{THME0064}. Rows show the observed target groups. For the
two other targets, only the search termination is shown.}
\label{tab:recorded-case}
\begin{tabularx}{\textwidth}{@{}llX@{}}
\toprule
Target group & Record ID & Evidence available before the decision \\
\midrule
\texttt{bravo}, \texttt{charlie} & & Payload search END at each target, zero exit status \\
\texttt{alpha} & 207481 & Payload search END, nonzero exit status \\
\texttt{alpha} & 207483 & Marker START \\
\texttt{alpha} & 207485 & Marker END, zero exit status \\
\bottomrule
\end{tabularx}
\end{table}
% 기존 원고의 target-wise GRU 대 relation attention 확률 비교 문장
% (0.434/0.450/0.481, 0.984/0.483/0.517)은 개요에 따라 부록 B.5로 옮긴다.
% 4회차 지시 9번: 표와 맞지 않던 "latest received result is therefore a success"를
% 고쳤다. alpha 그룹에서 마지막으로 받은 종료 상태는 marker END(0)이고 payload
% 검색 END는 nonzero다.
% 부록 B의 target-wise GRU 확률 비교는 한때 원 파일에서 찾지 못해 TODO로 바꿨으나, 2026-11-24
% sources/atlas_answers/defs.md 에서 원 보고서 값(.434/.450/.481)을 찾아 부록 B 에 되살렸다(방법 리뷰 P1-5).
% 이 장의 질문은 결론 뒤 "저자 확인 목록"에 모았다.

\section{What the errors look like}
\label{sec:diag-errors}

This appendix is a separate diagnosis of the in-house errors of
Chapter~\ref{ch:diagnosis}; CD does not depend on it. Every label of the 660 A+C rows is COPY or DISCOVERY, and every top-1
error of the 30 seeds on these rows, in both layouts, chose the other
of the two. For such a row the NLL splits exactly as
$-\log(p_{\mathrm{C}}+p_{\mathrm{D}})
 -\log\bigl(p_y/(p_{\mathrm{C}}+p_{\mathrm{D}})\bigr)$, where
$p_{\mathrm{C}}$ and $p_{\mathrm{D}}$ are the probabilities of COPY and
DISCOVERY and $y$ is the label. The first term, the probability missing
from the two file actions together, was small for every seed: at most
0.005 in either layout, and 0.002--0.005 on FLAT. What separates good
seeds from poor ones is the second term, the way a seed divides its
probability between COPY and DISCOVERY. This rules out the most direct
reading of the IP-hint inversion of Section~\ref{sec:conditions}, that a
seed takes the benign \texttt{ipconfig} records of worker~C for attacker
reconnaissance and predicts IP configuration enumeration. Whether the
benign IP hints shift the ratio between the two file actions indirectly
was not established.
% 근거: SEED_POPULATION FINAL_REPORT_KO.md:5 (분해 방식, ipconfig 직접 오인 설명
% 기각, 간접 효과 미확인), :37 ("30개 새 seed의 A+C 오답도 모두 파일 두 행동 사이",
% 총확률 부족 FLAT .0016--.0049, NESTED .0015--.0041).
% 660행 전부가 파일 행동: ORIG:PROPHET_MARKER_TARGET_RULE_20261119/REPORT_KO.md:7
% ("A+C FLAT/NESTED의 파일 행동 행 660개"), 341+319=660.
% 개요의 B′표에 따라 seed 181(이전 5-seed 집합)의 예는 쓰지 않았다.

On the seen condition, however, the problem is a different one. On
\texttt{ORIGINAL\_18} the seeds averaged 0.631, far above their A+C
losses, and most of this remaining loss reflects uncertainty about
which branch the procedure takes next rather than a weakness of the
encoder. We built a count table, the \emph{hint-branch table}, from
the 603 FIT runs. Its key is the receipt order of the hint-bearing START
records before the cutoff, together with whether the last of them has a
received END. For each key the table returns the smoothed action
frequencies observed under it in FIT. Used alone, the table scored 0.554 on the 118
\texttt{ORIGINAL\_18} rows and 0.496 on a separate, previously
inspected set of 64 runs and 472 rows of the seen condition. With each FIT
run held out in turn, it reached 82.5\% accuracy and an NLL of 0.376,
and splitting each key into two bins of the time since the last received
batch (age) lowered the NLL to 0.338. On the seen condition the
hint-branch table, which knows nothing about targets, thus did better
than the mean of the trained seeds (0.554 against 0.631). The trained
seeds were about 0.08 worse than a table that sees only the order of
the hints, so on this condition they have not used all that the hint
order offers; how much a better encoder could still gain there was not
measured.

Under the background workers the same table does not transfer. Benign START
records lengthen and reorder the hint sequences, to as many as 22 in
A+C against at most 15 in FIT, and for none of the 660 A+C rows does
the key occur in FIT; every A+C row falls back to the global action
frequency. The A+C losses must therefore be explained by something
other than knowledge of the branches.
% 근거: ORIG:PROPHET_HINT_BRANCH_TABLE_CORRECTED_20261118/REPORT_KO.md:7--11
% (키 정의, q_k 평활, 155개 FIT 키, 분기표 단독 0.553888 / 0.495788),
% :30--37 (FIT 실행 보류 NLL .375809 → 키+age 2구간 .337742, 정확도 .825180;
% age_s = 마지막 수신 batch 이후 cutoff까지 시간).
% ORIG:PROPHET_AC_HINT_BRANCH_COMPARE_V2_20261118/FINAL_REPORT_KO.md:3, 43
% (660 자연 경계 모두 FIT 키 지원 0 → 전역 빈도 π; FIT 최대 길이 15, A+C 22).
% RECORD:92--96. 472행 집합은 135실행 밖의 원래 과제 평가 차수다(장 첫 문단의 예외).
% 0.631은 같은 118행의 30 seed 평균(SEED_POPULATION FINAL_REPORT_KO.md:43).
% "공통 오답 13행" 판독은 개요 B′표에 따라 숫자 없이도 쓰지 않았다.
% "한 발 늦은 예측" 가설 기각은 11.4절에 숫자 없이 넣었다.

\section{A lookup rule without parameters}
\label{sec:diag-rule}

On the background executions the file actions follow a fixed order,
which we read off the row features of the development set. An
execution opens with four
payload searches (DISCOVERY). Each target is then visited with a
marker search (both terms are defined in Section~\ref{sec:recorded-case}), and COPY
follows if that target's payload search succeeded, that is, ended with
exit status zero; the source copies exactly when its observed result
for the current target is a success. The errors of REL\_GLOBAL on
COPY rows lie at the COPY decisions of the sixth, eighth and tenth
natural boundaries, and further errors fall on the boundary just after
a COPY, whose START has occurred but not yet been received.

This order suggests a rule with no
parameters. Let $T$ be the folder of the last received marker START.
Predict COPY if the last received payload END in $T$ has exit status
zero and no COPY START has been received since that marker; otherwise
predict DISCOVERY. On the boundaries where the next action depends on
the current target's result, the rule was correct on 620 of 620 FIT rows
and 476 of 476 background rows. Over all file-action boundaries it was
correct on 1,974 of 2,171 FIT rows and 1,290 of 1,309 background rows,
so it describes one decision of the procedure and is not an 18-way
predictor. Each of the 19 background errors predicted COPY again right
after an actual COPY whose START had not yet been received.

As a post hoc analysis on the same development set, we then turned the
rule into probabilities. Each rule category
received a COPY probability estimated from FIT runs, which replaced the
seed's split between COPY and DISCOVERY; the seed's total for the two
actions and its other 16 probabilities were kept. The mean single-seed
NLL fell from 0.142 to 0.091 on \texttt{A+C FLAT} and from 0.126 to
0.054 on \texttt{A+C NESTED}. It improved 19 of the 30 seeds on FLAT
and 26 on NESTED, and the mean number of wrong rows per seed fell from
17.4 to 7.4 and from 14.7 to 2.8. The same probabilities raised the loss
on \texttt{C\_ONLY} from 0.103 to 0.162, and on \texttt{B+C} from 0.059
to 0.062. Much of the A+C error thus sits in this one lookup, but a
probability calibrated on FIT runs does not carry over to every
background condition.
% 근거: RECORD:104--110. 절차 틀(처음 4개 DISCOVERY, marker 뒤 END ZERO면 COPY,
% source 정책 observed[target]==SUCCESS): RECORD:104--107.
% 출처(경계 6·8·10의 COPY 오답): atlas OT/PROPHET_MARKER_HEAD_INPUT_20261119/FINAL_REPORT_KO.md:42
%   ("경계 6·8·10의 정답 COPY 평균 오답 수는 FLAT에서 REL 8.47");
%   OT/PROPHET_MARKER_TARGET_RULE_20261119/FOLLOWUP_REPORT_KO.md:37 (diag.md 20a).
%   경계별 "손실 비중"표는 atlas에 없어 손실 집중이 아니라 오답 위치로 썼다.
% 규칙·620/620·476/476·1,974/2,171·1,290/1,309·배경 오답 19행(경계 7·9·11, native
% COPY START는 cutoff 전 발생, 미수신): ORIG:PROPHET_MARKER_TARGET_RULE_20261119/
% REPORT_KO.md:5, 13--23, 29--31.
% 확률화: FOLLOWUP_REPORT_KO.md:7--27 (범주별 q, 실행 보류, NO_MARKER는 seed 분할
% 보존, 파일 두 행동 총질량 m 유지; FLAT .141949→.091222, NESTED .125527→.054287,
% C_ONLY .103086→.162220, B+C .059181→.062372; 개선 seed 19/30, 26/30;
% 평균 오답 17.43→7.43, 14.70→2.83). ORIGINAL_18은 marker 0행이라 출력이 그대로다.
% 기록의 "seed 범위 .080~.140"(RECORD:69, 110)은 atlas에 없어 쓰지 않았다(diag.md 20b).

\section{Giving the model the lookup}
\label{sec:diag-intervention}

Inside the model, the lookup is not computed the same way across seeds.
This post hoc reading of the development set splits the pre-softmax
logits of each checkpoint into a base part and the relational residual. In the five seeds with the
highest A+C loss in each layout, 107 of the 148 errors on COPY rows
already had a residual that favored DISCOVERY, and in the other 41 the
base part overturned a residual that pointed the right way. The five
best seeds of each layout made 21 such errors.

We then gave the lookup's result to the model directly, first as a
development comparison on the same 135 runs. \texttt{OBSERVED\_HEAD} feeds
five indicators for the current target (payload END zero, payload END
nonzero, END not yet received, no marker, COPY START received after the
marker) through a 64-dimensional linear projection into the first
hidden layer of the output head, 320 parameters more than REL\_GLOBAL.
Its control, \texttt{SHUFFLED\_HEAD}, has the same width but takes the
indicators from another run with the same number of natural boundaries,
so the comparison separates the information from the extra input.
Table~\ref{tab:diag-rule} gives the results. On the development set,
\texttt{OBSERVED\_HEAD} lowered the A+C loss below both REL\_GLOBAL and
\texttt{SHUFFLED\_HEAD}.

The models were then frozen and checked on new runs. In the terms of
Section~\ref{sec:eval-rounds} this is a new-run check: it was sealed
before its runs were collected, but its 30 checkpoints had been trained
before it was written, so it is not a sealed criterion. On 157 runs
collected for it in the same lab, all four A+C contrasts against the two other models
had Bonferroni-adjusted 98.75\% intervals below zero, and the seed
standard deviation and the worst seed improved as well. These are the
157 runs that the 10-seed round of Chapter~\ref{ch:results} later
reused; by then they had been inspected. The gain held against both
models on a five-target FLAT variant. With three targets, the contrast
against \texttt{SHUFFLED\_HEAD} was not confirmed: on 32 runs and 256
rows of \texttt{A+C FLAT} the difference was $-0.019$ with a 98.75\%
interval of $[-0.043, +0.006]$. Preservation of the seen condition was not established. On 32 new
\texttt{ORIGINAL\_18} runs \texttt{OBSERVED\_HEAD} was 0.015 worse than
REL\_GLOBAL, inside the $+0.020$ limit, but the 95\% interval
$[-0.017, +0.048]$ crossed it, so the model was not adopted as a
general replacement.

We also built a generic version, \texttt{GENERIC\_HEAD}, and compared
it again on the 135 development runs. It describes the most recent
target key and its received results without the lab's marker and
payload commands. It did not keep the gain and did not beat its own
capacity-matched shuffled control.

Supplying the lab-specific lookup to the head thus helps, and the
lookup is a procedure-specific predictive bottleneck for this lab. It does not show
that the lookup causes the seed instability, because the features
encode the source policy directly, the generic form failed, and the
shuffled control does not preserve every temporal and structural
property of a run. We also did not find an architectural change built
on the lookup that is general and leaves the seen condition intact.
% 근거: RECORD:112--124.
% logit 분해(사후, 135실행): FOLLOWUP_REPORT_KO.md:31--44 (하위 5 seed COPY 오류
% 77+71=148 중 관계 부호 오류 56+51=107, 기본 경로 역전 21+20=41; 상위 5 seed 10+11=21).
% OBSERVED_HEAD 5특징·64차원·+320: RECORD:112; 개요 B표 일치.
% SHUFFLED_HEAD: 실행 간 derangement, 같은 자연 경계 수끼리 대응
% (ORIG:PROPHET_OBSERVED_HEAD_STAGE1_CONFIRM_20261119/FINAL_REPORT_KO.md:13).
% 1단계 새 실행: 같은 파일:3, 11, 19--37 (157실행·1,611행, 네 대비 Bonferroni 98.75%
% 상단 < 0, SD 차이 구간 < 0, 최악 seed FLAT .049529 대 REL .272791; ORIGINAL_18
% +0.015076, 95% [-0.017449, +0.047688]).
% 5대상: ORIG:PROPHET_OBSERVED_HEAD_STAGE2_FIVE_TARGET_FEASIBILITY_20261119/
% FINAL_REPORT_KO.md:3--15 (32실행·448행, 두 대비 보정 구간 상단 < 0).
% 3대상: 같은 파일:15 ("3대상에서 섞기 대비 구간이 0을 포함").
% 출처(3대상): atlas OT/PROPHET_OBSERVED_HEAD_STAGE2_3TARGET_20261119/FINAL_REPORT_KO.md:3,7-9,21
%   (OBSERVED−SHUFFLED −0.019314, 98.75% [−0.043177, +0.005801]; REL .171913, SHUFFLED .098368,
%   OBSERVED .079054; 원 97.5% 구간은 오류라 정정 구간을 따른다) (diag.md 20c).
% GENERIC_HEAD(135실행 개발 비교): ORIG:PROPHET_GENERIC_HEAD_CAPACITY_MATCHED_20261119/
% FINAL_REPORT_KO.md:3 (일반형−특화형 +0.053672/+0.049088, 98.75% 구간
% [+0.024236,+0.090602]/[+0.015770,+0.097816]; 일반형 섞기 대비 구간 0 포함;
% 201열 입력, 72,746 파라미터 대 특화형 60,202).
% 1단계 157실행 = 10-seed 차수의 평가 실행(PROPHET_ENV_CONSENSUS_KD_20261119/
% FINAL_REPORT_KO.md:12 "1단계 새 157실행은 이미 개봉한 개발 자료다").

\begin{table}[htbp]
\centering
\caption[Mean single-seed NLL of the lookup and of the head variants]{Mean single-seed NLL of the lookup and of the head variants
over the same 30 seeds. Development: the 135-run set, development
comparison; the lookup column is the post hoc probabilistic rule of
Section~\ref{sec:diag-rule} applied to REL\_GLOBAL. New runs: the
new-run check of Section~\ref{sec:diag-intervention}, collected in the
same lab with frozen checkpoints, and its five- and three-target FLAT
variants. On \texttt{ORIGINAL\_18} the lookup changes nothing, because
no row there has a marker.
``---'': not computed for that set.}
\label{tab:diag-rule}
\small
\setlength{\tabcolsep}{2.8pt}
\begin{tabular}{@{}llrrrrrr@{}}
\toprule
& & & REL\_GLOBAL & \multicolumn{2}{c}{Head with lookup features}
  & \multicolumn{2}{c}{Generic features} \\
\cmidrule(lr){5-6}\cmidrule(lr){7-8}
Set & Runs / rows & REL\_GLOBAL & + lookup (post hoc) & Shuffled & Observed
  & Shuffled & Generic \\
\midrule
\multicolumn{8}{@{}l}{\emph{Development set}} \\
\texttt{A+C FLAT}   & 31 / 341 & 0.142 & 0.091 & 0.106 & 0.065 & 0.106 & 0.119 \\
\texttt{A+C NESTED} & 29 / 319 & 0.126 & 0.054 & 0.091 & 0.059 & 0.093 & 0.108 \\
\texttt{C\_ONLY}    & 27 / 297 & 0.103 & 0.162 & 0.083 & 0.063 & 0.071 & 0.071 \\
\texttt{B+C}        & 32 / 352 & 0.059 & 0.062 & 0.041 & 0.024 & 0.050 & 0.037 \\
\texttt{ORIGINAL\_18} & 16 / 118 & 0.631 & 0.631 & 0.629 & 0.624 & 0.636 & 0.614 \\
\midrule
\multicolumn{8}{@{}l}{\emph{New runs, four targets}} \\
\texttt{A+C FLAT}     & 30 / 330 & 0.056 & --- & 0.054 & 0.016 & --- & --- \\
\texttt{A+C NESTED}   & 32 / 352 & 0.059 & --- & 0.044 & 0.019 & --- & --- \\
\texttt{C\_ONLY}      & 32 / 352 & 0.082 & --- & 0.064 & 0.045 & --- & --- \\
\texttt{B+C}          & 31 / 341 & 0.045 & --- & 0.036 & 0.024 & --- & --- \\
\texttt{ORIGINAL\_18} & 32 / 236 & 0.424 & --- & 0.433 & 0.439 & --- & --- \\
\midrule
\multicolumn{8}{@{}l}{\emph{New runs, five targets}} \\
\texttt{A+C FLAT}     & 32 / 448 & 0.060 & --- & 0.041 & 0.019 & --- & --- \\
\midrule
\multicolumn{8}{@{}l}{\emph{New runs, three targets}} \\
\texttt{A+C FLAT}     & 32 / 256 & 0.172 & --- & 0.098 & 0.079 & --- & --- \\
\bottomrule
\end{tabular}
\end{table}
% 근거: 개발 집합 REL_GLOBAL·섞기·OBSERVED·일반형 섞기·GENERIC:
% GENERIC_HEAD_CAPACITY_MATCHED FINAL_REPORT_KO.md:17--22 (.141949/.106273/.065357/
% .105593/.119028; .125527/.091312/.058858/.092673/.107946; .103086/.082897/.063363/
% .071399/.071087; .059181/.040843/.023705/.049861/.037245).
% + lookup 열: FOLLOWUP_REPORT_KO.md:17--22 (단일 seed 평균, 고정 q 적용).
% 새 실행 4대상: OBSERVED_HEAD_STAGE1_CONFIRM FINAL_REPORT_KO.md:19--25.
% 새 실행 5대상: OBSERVED_HEAD_STAGE2_FIVE_TARGET FINAL_REPORT_KO.md:5--9.
% 개발 집합과 새 실행은 다른 평가 실행이므로 칸끼리 빼거나 합치지 않는다.
% 3대상 행: OBSERVED_HEAD_STAGE2_3TARGET FINAL_REPORT_KO.md:7-9.
% 개발 ORIGINAL_18 행: atlas OT/PROPHET_MARKER_HEAD_INPUT_20261119/FINAL_REPORT_KO.md:27;
%   OT/PROPHET_GENERIC_HEAD_CAPACITY_MATCHED_20261119/FINAL_REPORT_KO.md:21 (.630865/.629411/.623849/
%   .636115/.613693). lookup 열은 marker 0행이라 REL과 같다(FOLLOWUP:29,31) (diag.md 20d).

% 이 장의 질문은 결론 뒤 "저자 확인 목록"에 모았다.
% 6.3--6.5절 위치는 지금 자리에 둔다(sources/atlas_answers/decisions.md 22번).
